\subsection{可一致化空间上的半连续函数}

在第四章，§ 6，第 2 号，定理 4 的推论中，我们证明了：拓扑空间上一族连续实值函数的上包络是下半连续函数。若空间可一致化，这个命题有一个逆命题：

命题 5. 为使每个下半连续实值函数 \(f\) （有限或否，定义在拓扑空间 \(\mathbf{X}\) 上）都是 \(\mathbf{X}\) 上的连续实值函数（有限或否）之 \(\leqslant f\) 者的上包络，必要且充分的是 \(\mathbf{X}\) 可一致化。

该条件是必要的。设 \(x_0\) 是 \(\mathbf{X}\) 中任意一点，\(\mathbf{V}\) 是一个开邻域（ \(x_0\) 的）；则特征函数 \(\varphi_{\mathbf{V}}\) —— 集 \(\mathbf{V}\) 的特征函数 —— 是下半连续的（第四章，§ 6，第 2 号，命题 1 的推论）；于是由假设，存在一个连续实值函数 \(g\) （在 \(\mathbf{X}\) 上），使得 \(g \leqslant \varphi_{\mathbf{V}}\) 且 \(g(x_0) = a > 0\) 。连续函数 \(\inf\left(\mathbf{i}, \frac{\mathbf{i}}{a} g^+\right)\) 取值于 {[}o, i{]}，在 CV 中等于 o ，在 \(x_0\) 处等于 i 。因而（定理 2）\(\mathbf{X}\) 可一致化。

该条件是充分的。先考察 \(f\) 取值于 \([-1, +1]\) 的情形。我们要证明：对每个 \(x_0 \in X\) 与每个数 \(a < f(x_0)\) ，存在一个连续实值函数 \(g\) （在 \(X\) 上），使得 \(g \leqslant f\) 且 \(g(x_0) \geqslant a\) 。若 \(a \leqslant -1\) ，可取 \(g\) 为常数 \(-1\) 。若 \(-1 < a < f(x_0)\) ，则存在一个邻域 \(V\) （ \(x_0\) 的），使得 \(f(x) \geqslant a\) 对一切 \(x \in V\) 成立。由于 \(X\) 可一致化，存在一个连续实值函数 \(h\) （在 \(X\) 上，取值于 \([0, 1]\) ），使得 \(h(x_0) = 0\) 且 \(h(x) = 1\) 对一切 \(x \in C_V\) 成立。于是可取 \(g(x) = a - (a + 1)h(x)\) ，我们就得到一个满足所述条件的连续函数。注意这个函数取值于 \([-1, +1]\) 。

一般情形由结构转移得到；因为存在 \([-1, +1]\) 到 \(\overline{R}\) 上的严格递增同胚（第四章，§ 4，第 2 号，命题 2），而半连续函数的定义只涉及 \(\overline{R}\) 的序结构与拓扑。

注。在上面的证明中我们看到，函数 g 不取值 +1。由结构转移可得：可一致化空间 X 上的每个下半连续实值函数 f 都是 X 上不取值 +∞ 的连续实值函数 \(g \leqslant f\) 的上包络。

